% RL-CONIC-2026 / theoretical / T4 -> T5 provisional / MW-MASTER-V04
% Author name is exactly 菠萝, as explicitly required by the author.
% No affiliation, email, or funding is invented.
\documentclass[11pt]{article}
\usepackage[T1]{fontenc}
\usepackage[utf8]{inputenc}
\usepackage{lmodern}
\usepackage[a4paper,margin=28mm]{geometry}
\usepackage{amsmath,amssymb,amsthm,mathtools}
\usepackage{enumitem}
\usepackage{booktabs,array}
\usepackage{xcolor}
\usepackage{microtype}
\usepackage[numbers,sort&compress]{natbib}
\usepackage[unicode=true,colorlinks=true,linkcolor=blue!45!black,citecolor=blue!45!black,
  urlcolor=blue!45!black,breaklinks=true]{hyperref}
\hypersetup{pdftitle={Frozen minimal-face constant rank and metric subregularity over amenable cones},
  pdfauthor={菠萝},pdfsubject={Theoretical manuscript v0.4}}
% Native pdfLaTeX font embedding for the two Unicode author-name glyphs.
\pdfmapline{+author-cjk AuthorNameCJK "AuthorNameCJKEncoding ReEncodeFont" <author-cjk.enc <author-cjk.ttf}
\pdfglyphtounicode{uni83E0}{83E0}
\pdfglyphtounicode{uni841D}{841D}
\newcommand{\authorname}{\begingroup
  \font\authorcjk=author-cjk at 12pt\authorcjk\char0\char1\endgroup}
\setlength{\emergencystretch}{2em}
\setlist{itemsep=3pt,topsep=5pt}
\numberwithin{equation}{section}
\newtheorem{theorem}{Theorem}[section]
\newtheorem{lemma}[theorem]{Lemma}
\newtheorem{proposition}[theorem]{Proposition}
\newtheorem{corollary}[theorem]{Corollary}
\theoremstyle{definition}
\newtheorem{definition}[theorem]{Definition}
\newtheorem{example}[theorem]{Example}
\theoremstyle{remark}
\newtheorem{remark}[theorem]{Remark}
\DeclareMathOperator{\dist}{dist}
\DeclareMathOperator{\ri}{ri}
\DeclareMathOperator{\rbd}{rbd}
\DeclareMathOperator{\spann}{span}
\DeclareMathOperator{\range}{range}
\DeclareMathOperator{\rank}{rank}
\DeclareMathOperator{\cone}{cone}
\DeclareMathOperator{\conv}{conv}
\DeclareMathOperator{\tr}{tr}
\DeclareMathOperator{\diag}{diag}
\DeclareMathOperator{\cl}{cl}
\DeclareMathOperator{\interior}{int}
\newcommand{\R}{\mathbb R}
\newcommand{\Ssym}{\mathbb S}
\newcommand{\xb}{\bar x}
\newcommand{\xhat}{\widehat x}
\newcommand{\ip}[2]{\langle #1,#2\rangle}
\newcommand{\norm}[1]{\lVert #1\rVert}
\newcommand{\dset}[2]{\dist(#1,#2)}

\title{Frozen minimal-face constant rank and metric subregularity\protect\\
  over amenable cones}
\author{\authorname}
\date{September 9, 2026\\[3pt]\small Working manuscript v0.4}

\begin{document}
\maketitle

\begin{abstract}
We study continuously differentiable conic systems at a cone apex under a
constant-rank condition defined by the reference minimal linearized face.
For a finite-dimensional proper nice cone, reference closedness of the
adjoint cone image and constancy of the normal rank imply closedness of
nearby adjoint images, stability of the minimal linearized face, and local
nonlinear facial reduction. The key observation is a rank sandwich on two
fixed dual-face subspaces. A quantitative normal correction then moves
each nearby input to one common face manifold with displacement controlled
by the original cone residual. If the reference face is amenable in the
cone, a second correction within that manifold yields metric
subregularity and an explicit local upper bound for its modulus.
Consequently, within the class of proper nice cones, amenability is
equivalent to the universal validity of this constant-rank-to-subregularity
implication for all apex systems. The reverse implication is already
detected by linear face inclusions. An adaptation of the published
four-dimensional nice nonamenable cone of Louren\c{c}o, Roshchina, and
Saunderson shows failure of the error bound even under full facial CRCQ.
A coupled second-order-cone system and positive-semidefinite systems with
high-dimensional nonpolyhedral faces illustrate the joint-input argument.
\end{abstract}

\noindent\textbf{Keywords.} Constant rank; constraint qualification;
minimal face; nice cone; amenable cone; metric subregularity; facial reduction.

\section{Introduction}
\label{sec:introduction}

For a smooth constraint mapping $G$ and a closed convex cone $C$, a local
error bound compares the distance of an input to $G^{-1}(C)$ with the
output residual $\dset{G(x)}{C}$. Such a comparison is particularly useful
when a strict feasible linearized direction is unavailable. A face-based
constant-rank condition can describe this degeneracy without requiring
the original linearized system to meet the interior of the whole cone.
The resulting question has two distinct parts: does the condition
stabilize the appropriate face, and does it control distance using the
\emph{original} cone residual?

These questions should not be conflated. Local equality between the
original feasible set and a face-reduced feasible set is a geometric
statement. It does not automatically compare the distances to the cone
and to the face. The present paper separates these roles. Niceness
provides the closed-image geometry needed for face stability and normal
correction. Amenability of the particular reference face supplies the
metric comparison that completes the error bound.

\newpage
\paragraph{Main results.}
For $C^1$ apex systems with the frozen-face condition defined below,
the results have three parts.
\begin{enumerate}[label=\arabic*.]
\item Over a proper nice cone, frozen minimal-face CRSC gives nearby
closed adjoint images and the same minimal linearized face
(Theorem~\ref{thm:stability}), together with local nonlinear facial
reduction (Corollary~\ref{cor:fr}).
\item If the reference face is amenable in the original cone, the same
condition yields MSCQ with an explicit sufficient local constant
(Theorem~\ref{thm:mscq}); amenable cones satisfy this face hypothesis
(Corollary~\ref{cor:amenable}).
\item For a fixed proper nice cone, amenability is equivalent to the
universal CRSC-to-MSCQ property over all admissible apex systems
(Theorem~\ref{thm:universal}); the nice nonamenable example fails MSCQ
even under full facial CRCQ (Proposition~\ref{prop:counterexample}).
\end{enumerate}

\subsection{Relation to the closest results}

The minimal-face CRSC framework of Andreani, Haeser, Mito, and
Ram\'irez \citep{ahmr2026} is the nearest starting point. Our numbered
comparison refers specifically to their March 20, 2025 author manuscript
\citep{ahmr2025author}. Its Definition~3.2 freezes the minimal linearized
face at the reference point. Proposition~5.2 proves nonlinear facial
reduction from a strict conjugate-face kernel multiplier and a
constant-rank condition on the conjugate-face span. Theorem~5.2 obtains
the reduction with additional nearby closed-image and face-stability
assumptions, denoted A1 and A2, and the following paragraph asks whether
they can be removed. We derive the missing conjugate-span rank condition
directly from frozen minimal-face CRSC, and also prove A1 and A2
themselves on an open input neighborhood. This removes the additional
assumptions from the theorem as stated in that author version. The
comparison is not a claim that the final published text has been checked
to be identical to that version.

Our geometric argument imports Pataki's criterion for a closed linear
image of a cone \citep[Theorem~1.1]{pataki2007}. Its exact equality of a conjugate-face image with a
linear subspace is essential. Replacing this equality by an equality
only after taking closures would not justify the rank argument below.
The other imported ingredient is cone amenability: a facewise error
bound on the linear span of a face. This concept, its relation to
niceness, and its product and standard-cone properties are established
results of Louren\c{c}o \citep{lourencoAmenable} and
Louren\c{c}o, Roshchina, and Saunderson \citep{lrs2022}.
Those geometric properties are not claimed as contributions here.
Nor is the short rank sandwich a new general principle of facial
reduction: it supplies the hypothesis needed by an existing proposition.

The nonlinear step is necessary because a general value $G(x)$ need not
belong to the span of the minimal face, where amenability applies. We
first produce an input correction to that span with cost measured by the
original residual. Only then do we invoke amenability and perform a
second, relative-interior correction. Both steps use the same joint input
mapping. No intersection regularity is inferred from separate blockwise
error bounds.

Related error-bound results have different model hypotheses.
\citet[Theorem~5.1]{chieu2026soc} establish a CRCQ--MSCQ equivalence for
an affine single-SOC system. The nonlinear theorem of
\citet[Theorem~4.1]{huy2025cone} uses $C^1$ cone-concave data, a specified
smooth-cone class, and a nearby Abadie condition. These hypotheses are
not assumed here. Table~\ref{tab:nearest} records the distinctions that
are relevant to the proof, without implying that all other error-bound
results have been excluded.

\begin{table}[!htbp]
\centering\small
\caption{Closest interfaces and their relevant hypotheses. The first row
uses the explicitly cited author version.}
\label{tab:nearest}
\begin{tabular}{@{}>{\raggedright\arraybackslash}p{.24\linewidth}
                  >{\raggedright\arraybackslash}p{.34\linewidth}
                  >{\raggedright\arraybackslash}p{.36\linewidth}@{}}
\toprule
Source & Relevant statement & Interface with this paper\\
\midrule
\citet{ahmr2025author}, Theorem~5.2 & Frozen minimal-face CRSC with A1/A2;
nonlinear facial reduction & The fixed-subspace rank argument supplies
the missing rank condition and yields A1/A2.\\[4pt]
\citet{lourencoAmenable}; \citet{lrs2022} & Face-to-cone residual comparison
on $\spann F$ & Imported after a quantitative input correction to the
common face manifold.\\[4pt]
\citet{chieu2026soc}, Theorem~5.1 & Affine single-SOC CRCQ and MSCQ &
The present sufficient theorem allows arbitrary $C^1$ joint maps and
amenable output cones.\\[4pt]
\citet{huy2025cone}, Theorem~4.1 & Smooth cone, cone-concavity, and nearby
Abadie condition & The present hypothesis is frozen minimal-face CRSC;
cone-concavity is not imposed.\\
\bottomrule
\end{tabular}
\end{table}

Facial residual functions allow non-Lipschitz moduli for affine conic
feasibility \citep{lindstrom2023frf}; radial-type estimates for
semidefinite feasibility likewise concern an affine-space/PSD model
\citep{nishijima2026radial}. Such results do not supply, merely by their
generality in the output cone, the nonlinear input correction used here.
Modern facial-reduction and joint-supporting-subspace approaches for
convex programs \citep{lin2025facial,scott2025joint} have convex-set or
$f+g\circ A$ models; an arbitrary $C^1$ inverse image need not be convex.
Recent constant-rank-type local error bounds and non-relaxable-set
qualifications \citep{andreani2026new,andreani2026nonrelaxable} concern
scalar nonlinear programming. In particular, the Lower-CRSC testing
domain in the latter is not the open-neighborhood frozen-face condition
used below. These comparisons locate the present hypotheses; they do
not claim a first constant-rank error bound in nonlinear optimization.

\subsection{Organization}

Section~\ref{sec:preliminaries} fixes the notation and condition.
Sections~\ref{sec:nice} and \ref{sec:mscq} develop the two corrections
and the universal characterization. Section~\ref{sec:boundary} gives
the apex counterexample. Section~\ref{sec:examples} treats cone classes,
two coupled nonlinear examples, and fixed local reductions.
Section~\ref{sec:residual-interface} gives the regularity interface
through a residual comparison.

\section{Preliminaries and the frozen-face condition}
\label{sec:preliminaries}

All spaces are finite-dimensional Euclidean spaces and all distances are
Euclidean unless a subspace metric is explicitly indicated. For a linear
map $A:X\to E$, its adjoint is $A^*:E\to X$. For a subspace $L$, $P_L$
denotes the orthogonal projection. We use the positive dual
\[
 C^*=\{y\in E:\ip{y}{c}\geq0\text{ for every }c\in C\}.
\]
The negative-polar convention changes only the sign of an adjoint cone
image, not its closedness or rank.

A cone is \emph{proper} here if it is closed, convex, pointed, and
full-dimensional in a positive-dimensional ambient space. If $F$ is a
face of $C$, then
\[
 F^\perp=(\spann F)^\perp,
 \qquad F^\triangle=C^*\cap F^\perp.
\]
We use $F_{\min}(Q)$ for the smallest face of $C$ containing $Q\subset C$.
A convex set meets the relative interior of its minimal containing face.
In particular, if $Q=\range A\cap C$, there is $\bar d\in X$ such that
$A\bar d\in\ri F_{\min}(Q)$. For a nonzero face the image can be chosen
nonzero. One may see this by taking finitely many points whose faces
generate the minimal face and summing them. The set $\range A\cap C$ is
a convex cone, so the sum remains in it.

\begin{definition}[Niceness and face amenability]
\label{def:amenable}
A closed convex cone $C$ is \emph{nice} if $C^*+J^\perp$ is closed for
every face $J$ of $C$. A face $J$ is \emph{amenable in $C$} if there is a
finite constant $a_J>0$ such that
\begin{equation}
 \dset{y}{J}\leq a_J\dset{y}{C}
 \qquad (y\in\spann J).
 \label{eq:amenable}
\end{equation}
The cone is \emph{amenable} if every face is amenable in this sense.
\end{definition}

Amenable cones are nice, and nice cones are facially exposed. These are
imported cone-geometric facts \citep{lourencoAmenable,lrs2022}. In
particular, for a nice cone and a proper face $F$, a point in
$\ri F^\triangle$ exposes $F$; the conjugate face is nonzero. No
constant uniform over all faces is required in
Definition~\ref{def:amenable}.

Fix an open set $V\subset X$, a point $\xb\in V$, and a $C^1$ map
$G:V\to E$ with $G(\xb)=0$. Write
\begin{equation}
 \Omega=\{x\in V:G(x)\in C\},\qquad A=DG(\xb),\qquad
 F=F_{\min}(\range A\cap C),
 \label{eq:setup}
\end{equation}
and freeze the subspaces
\begin{equation}
 H=F^\perp,\qquad S=\spann F^\triangle\subset H.
 \label{eq:subspaces}
\end{equation}
Every local assertion below is understood after shrinking to a
neighborhood of $\xb$ contained in $V$.

\begin{definition}[Frozen minimal-face CRSC]
\label{def:crsc}
The system in \eqref{eq:setup} satisfies frozen minimal-face CRSC at
$\xb$ if $A^*C^*$ is closed and there is an open neighborhood of $\xb$
on which
\begin{equation}
 \rank\bigl(DG(x)^*|_H\bigr)=r,
 \qquad r=\rank(A^*|_H).
 \label{eq:crsc}
\end{equation}
\end{definition}

\begin{definition}[Full facial CRCQ at an apex]
\label{def:crcq}
The facial constant-rank property holds at $\xb$ if there is one
neighborhood $W$ of $\xb$ such that, for every face $J$ of $C$,
$\rank(DG(x)^*|_{J^\perp})$ is constant for $x\in W$; the constant
may depend on $J$. Full facial CRCQ consists of this property and
closedness of $A^*C^*$.
\end{definition}
This is Definition~3.1 of \citet{ahmr2025author} in the apex
representation with the identity reduction. The faces tested are all
faces of $C$, not only subfaces of $F$. Full facial CRCQ implies
Definition~\ref{def:crsc}.

The face $F$ and the reduction $G$, if a reduction is used, are fixed at
$\xb$. The rank condition is imposed on an open input neighborhood, not
only on feasible points. Nearby closedness of $DG(x)^*C^*$ and nearby
constancy of the minimal linearized face are \emph{conclusions}, not
additional assumptions.

\begin{definition}[MSCQ]
The system satisfies the metric subregularity constraint qualification
(MSCQ) at $\xb$ if some $\kappa>0$ and some neighborhood $W$ of $\xb$
satisfy
\begin{equation}
 \dset{x}{\Omega}\leq\kappa\dset{G(x)}{C}
 \qquad (x\in W).
 \label{eq:mscq}
\end{equation}
\end{definition}

For clarity, the closed-image theorem that is imported throughout is
stated explicitly. It is not proved or claimed as new here.

\begin{theorem}[Pataki's closed-image criterion, imported
{\citep[Theorem~1.1]{pataki2007}}]
\label{thm:pataki}
Let $C$ be proper and nice, let $A:X\to E$ be linear, and define $F,H,S$
as in \eqref{eq:setup}--\eqref{eq:subspaces}. The following are equivalent:
\begin{enumerate}[label=(\roman*)]
 \item $A^*C^*$ is closed;
 \item $A^*F^\triangle=A^*H$;
 \item both of the following hold:
 \begin{equation}
 \ker A^*\cap\ri F^\triangle\ne\varnothing,
 \qquad
 \range A\cap S^\perp=\range A\cap\spann F.
 \label{eq:pataki-kernel}
 \end{equation}
\end{enumerate}
\end{theorem}

\begin{remark}[Zero-dimensional reduction]
If a fixed local reduction has zero-dimensional output space, the reduced
map and cone are both zero. Local equivalence then makes every sufficiently
near input feasible, and any positive error-bound constant is valid.
This branch is separate from the positive-dimensional proper-cone setup;
no unit direction, positive singular value, or relative-boundary margin
will be invoked for it.
\end{remark}

\section{Nice cones: nearby stability and normal correction}
\label{sec:nice}

Throughout this section, $C$ is proper and nice and frozen minimal-face
CRSC holds at $\xb$.

\subsection{A rank sandwich on fixed subspaces}

\begin{lemma}[Fixed-subspace rank sandwich]
\label{lem:sandwich}
On some open neighborhood of $\xb$,
\begin{equation}
 \rank\bigl(DG(x)^*|_S\bigr)
 =\rank\bigl(DG(x)^*|_H\bigr)=r,
 \qquad DG(x)^*S=DG(x)^*H.
 \label{eq:sandwich}
\end{equation}
\end{lemma}
\begin{proof}
Theorem~\ref{thm:pataki} and reference closedness give
$A^*F^\triangle=A^*H$. Taking linear spans yields $A^*S=A^*H$,
so the rank of $A^*|_S$ is $r$. If $r>0$, choose fixed orthonormal bases
and a nonzero $r$-by-$r$ minor of this restricted matrix. Continuity of
$DG$ preserves that minor on a neighborhood. On the other hand,
$S\subset H$ and \eqref{eq:crsc} give
\[
 r\leq\rank(DG(x)^*|_S)
 \leq\rank(DG(x)^*|_H)=r.
\]
For $r=0$ the conclusion follows from the second inequality alone.
The subspace images are nested and have equal dimension, hence are equal.
\end{proof}

\begin{lemma}[Continuity of kernel projections]
\label{lem:kernel}
If a continuous matrix field $B(x)$ has constant rank near $\xb$, the
orthogonal projection onto $\ker B(x)$ depends continuously on $x$.
\end{lemma}
\begin{proof}
Choose a nonzero maximal minor at $\xb$ and retain it on a smaller
neighborhood. Solving the corresponding dependent coordinates expresses
the kernel as a graph of a continuous linear-map field over a fixed
coordinate subspace. If $Q(x)$ is the continuous full-column-rank matrix
whose columns span this graph, its orthogonal projection is
$Q(x)(Q(x)^*Q(x))^{-1}Q(x)^*$. The formula is continuous.
The zero-dimensional and full-dimensional kernels are included directly.
\end{proof}

\begin{theorem}[Full-neighborhood stability]
\label{thm:stability}
There is an open neighborhood $W$ of $\xb$ such that, for every $x\in W$,
\begin{equation}
 DG(x)^*C^*\text{ is closed},
 \qquad F_{\min}(\range DG(x)\cap C)=F.
 \label{eq:stability}
\end{equation}
Thus the nearby conditions A1 and A2 are automatic under the frozen
condition in Definition~\ref{def:crsc}.
\end{theorem}
\begin{proof}
Write $A_x=DG(x)$. First suppose $F=C$. The minimal-face property gives
a vector $\bar d$ with $A\bar d\in\interior C$. By continuity,
$A_x\bar d\in\interior C$ nearby, so the minimal face remains $C$.
Choose a common $\delta>0$ with $B(A_x\bar d,\delta)\subset C$
nearby. For $y\in C^*$,
\[
 \delta\norm{y}\leq\ip{A_x\bar d}{y}
 =\ip{\bar d}{A_x^*y}.
\]
For each fixed nearby $x$, every convergent sequence $A_x^*y_j$,
$y_j\in C^*$, therefore has bounded lifts $y_j$. A convergent
subsequence of the lifts proves that $A_x^*C^*$ is closed.
No surjectivity of $A_x$ is needed.

Suppose now that $F$ is proper. Choose
$Y_0\in\ri F^\triangle\cap\ker A^*$ using
Theorem~\ref{thm:pataki}. Lemmas~\ref{lem:sandwich} and
\ref{lem:kernel} show that
\[
 Y_x=P_{\ker(A_x^*|_S)}Y_0
 \longrightarrow Y_0,
\]
where the kernel and projection are taken inside $S$.
Since $\ri F^\triangle$ is open relative to $S$, shrinking the
neighborhood gives
\begin{equation}
 Y_x\in\ri F^\triangle,\qquad A_x^*Y_x=0.
 \label{eq:exposing}
\end{equation}
For $v\in\range A_x\cap C$, we have $\ip{Y_x}{v}=0$.
Indeed, for each $y\in F^\triangle$, relative interior gives
$Y_x-\varepsilon y\in F^\triangle$ for some $\varepsilon>0$.
Nonnegativity of both pairings with $v$ then forces $\ip{y}{v}=0$.
Facial exposedness, equivalently the double-conjugate identity for $F$,
therefore implies
\begin{equation}
 \range A_x\cap C\subset F.
 \label{eq:face-inclusion}
\end{equation}
If $F=\{0\}$ this already proves minimal-face stability.

If $F\ne\{0\}$, select $\bar d$ with $A\bar d\in\ri F$ and define
\[
 d_x=P_{\ker(P_HA_x)}\bar d.
\]
The map $P_HA_x$ has rank $r$, and $P_HA\bar d=0$.
Lemma~\ref{lem:kernel} gives $d_x\to\bar d$. Consequently
$A_xd_x\in\spann F$ and $A_xd_x\to A\bar d\in\ri F$.
The image is in $\ri F$ nearby. Together with
\eqref{eq:face-inclusion}, this proves the second assertion in
\eqref{eq:stability}.

Finally, the image equality in \eqref{eq:sandwich} gives
\[
 \ker(P_SA_x)=\ker(P_HA_x),
 \qquad
 \range A_x\cap S^\perp=\range A_x\cap\spann F.
\]
The second equality follows by representing a vector in the range as
$A_xd$ and applying the kernel equality. We have already proved that its
minimal face at $x$ is $F$. This equality and \eqref{eq:exposing} now
apply Theorem~\ref{thm:pataki} at $x$, proving closedness.
There is no use of nearby closedness earlier in the proof.
\end{proof}

\subsection{A normal angle and a common face manifold}

Define
\begin{equation}
 U=\range(P_HA),\qquad D=\cl(P_HC)\subset H,
 \qquad \phi=P_HG.
 \label{eq:normal}
\end{equation}
The closure in the definition of $D$ is deliberate: a projected
nonpolyhedral cone need not be closed.

\begin{lemma}[Positive normal angle]
\label{lem:angle}
The subspace $U$ and the closed cone $D$ satisfy $U\cap D=\{0\}$.
If $r>0$, then
\begin{equation}
 \eta=\min_{\substack{u\in U\\\norm{u}=1}}\dset{u}{D}>0.
 \label{eq:eta}
\end{equation}
\end{lemma}
\begin{proof}
Let $u=P_HAd\in D$. Every $y\in F^\triangle$ has
$\ip{u}{y}\geq0$, first on $P_HC$ and then by continuity on its closure.
For any $h\in H$, Theorem~\ref{thm:pataki} supplies
$y\in F^\triangle$ with $A^*y=A^*h$. Since both $h$ and $y$ are in $H$,
\[
 \ip{u}{h}=\ip{Ad}{h}=\ip{Ad}{y}=\ip{u}{y}\geq0.
\]
Applying the same argument to $-h$ gives $\ip{u}{h}=0$ for every
$h\in H$. Since $u\in H$, it is zero. If $r>0$, compactness of the unit
sphere in $U$ and closedness of $D$ give \eqref{eq:eta}.
\end{proof}

\begin{lemma}[Quantitative correction to the common manifold]
\label{lem:normal-correction}
Let $M=\{x\in V:\phi(x)=0\}$. Every sufficiently near $x$ admits
$\xhat\in M$ such that
\begin{equation}
 \norm{x-\xhat}\leq b\dset{G(x)}{C}.
 \label{eq:normal-correction}
\end{equation}
For $r>0$ one may take
\begin{equation}
 b=\frac{4}{\sigma\eta},\qquad
 \sigma=\sigma_{\min}^{+}(P_HA),
 \label{eq:b}
\end{equation}
where $\sigma_{\min}^{+}$ is the smallest positive singular value.
For $r=0$, one may take $b=0$ and $\xhat=x$.
\end{lemma}
\begin{proof}
The derivative $D\phi(x)=P_HDG(x)$ has rank $r$ near $\xb$.
If $r=0$, $\phi$ is constant on a small connected neighborhood, and is
zero because $\phi(\xb)=0$. This proves the stated branch.

Let $r>0$, put $N=\ker D\phi(\xb)$, and consider
\[
 \Theta(x)=\bigl(P_U\phi(x),P_N(x-\xb)\bigr)\in U\times N.
\]
Its derivative at $\xb$ is invertible: $D\phi(\xb)$ is an isomorphism
from $N^\perp$ to $U$, while the second component is the identity on $N$.
The inverse function theorem provides a $C^1$ inverse $x=\theta(u,z)$
on a small product neighborhood of $(0,0)$.

The derivative of $P_U\phi$ has rank $r$ near $\xb$, as does
$D\phi$. Its kernel therefore equals the kernel of $D\phi$.
Along each connected $z$-fiber the full map $\phi\circ\theta$ has zero
$z$-derivative. Since its $U$-component is $u$, it follows that
\begin{equation}
 \phi(\theta(u,z))=u+\beta(u),\qquad
 \beta(u)\in U^\perp\cap H,\quad \beta(0)=D\beta(0)=0.
 \label{eq:graph}
\end{equation}
This is precisely the constant-rank consequence needed here; it uses
only $C^1$ regularity.

At the reference point, $\norm{D_u\theta}=1/\sigma$.
Shrink the product chart so that
\[
 \norm{D_u\theta}\leq\frac{2}{\sigma},
 \qquad \norm{D\beta}\leq\frac{\eta}{2}.
\]
Cone homogeneity, \eqref{eq:eta}, and the distance Lipschitz inequality give
\[
 \dset{\phi(x)}{D}
 \geq\dset{u}{D}-\norm{\beta(u)}
 \geq\frac{\eta}{2}\norm{u}.
\]
Take $\xhat=\theta(0,z)$. Equation~\eqref{eq:graph} implies
$\xhat\in M$, and the segment in the $u$-coordinate yields
\[
 \norm{x-\xhat}\leq\frac{2}{\sigma}\norm{u}.
\]
Since $P_H$ is nonexpansive and $D$ contains $P_HC$,
$\dset{\phi(x)}{D}\leq\dset{G(x)}{C}$.
Combining these estimates proves \eqref{eq:normal-correction}.
Nested neighborhoods ensure the input and its correction remain in the
same chart.
\end{proof}

\begin{corollary}[Nonlinear facial reduction]
\label{cor:fr}
On a neighborhood of $\xb$,
\begin{equation}
 G^{-1}(C)=G^{-1}(F).
 \label{eq:fr}
\end{equation}
\end{corollary}
\begin{proof}
If $G(x)\in C$, the residual on the right of
\eqref{eq:normal-correction} is zero; its correction thus equals $x$.
Hence $\phi(x)=0$ and $G(x)\in C\cap\spann F=F$.
The reverse inclusion is immediate.
\end{proof}

\begin{remark}[Dependence on the existing facial-reduction argument]
The rank sandwich also supplies the conjugate-span rank hypothesis of
\citet[Proposition~5.2]{ahmr2025author}, so
Corollary~\ref{cor:fr} could be obtained from that result.
Lemma~\ref{lem:normal-correction} gives a direct proof and the additional
quantitative correction needed for the original-residual estimate below.
The full-face branch $F=C$ has $H=\{0\}$ and $M=V$ locally; no
nonzero conjugate-face multiplier is asserted in that branch.
\end{remark}

\section{Amenability and the original-residual error bound}
\label{sec:mscq}

\subsection{A second correction on the same manifold}

\begin{lemma}[Relative-interior correction]
\label{lem:relative}
Under the standing hypotheses of Section~\ref{sec:nice}, including
properness and niceness of $C$ and frozen minimal-face CRSC at $\xb$,
let $M$ be the common manifold in Lemma~\ref{lem:normal-correction}.
Assume $F\ne\{0\}$. There is a unit vector
$d_0\in\ker(P_HA)$ such that $v_0=Ad_0\in\ri F$. Define
\begin{equation}
 \tau_F=\dist_{\spann F}(v_0,\rbd F)>0.
 \label{eq:tau}
\end{equation}
For every sufficiently near $\xhat\in M$, there is $x^+\in M$ with
$G(x^+)\in F$ and
\begin{equation}
 \norm{x^+-\xhat}\leq c\dset{G(\xhat)}{F},
 \qquad c=\frac{4}{\tau_F}.
 \label{eq:relative}
\end{equation}
\end{lemma}
\begin{proof}
The minimal-face property gives a direction with nonzero image in
$\ri F$. Rescale it to obtain $d_0$; positive rescaling preserves the
relative interior. Since $Ad_0\in\spann F$, $d_0$ belongs to
$\ker(P_HA)$. The margin \eqref{eq:tau} is finite and positive because
$F$ is a nonzero pointed cone in its span.

Use local coordinates $m(z)$ for $M$ with $m(0)=\xb$, choosing
$m(z)=\theta(0,z)$ in the positive-rank construction and the identity
chart in the rank-zero case. There is a fixed coordinate direction $w_0$
with $Dm(0)w_0=d_0$. The map $G\circ m$ takes values in $\spann F$.
By continuity, shrink the chart so that
\[
 \norm{Dm(z)w_0}\leq2,
 \qquad
 \norm{D(G\circ m)(z)w_0-v_0}\leq\frac{\tau_F}{4}.
\]
For $\xhat=m(z)$ set $q=\dset{G(\xhat)}{F}$. If $q=0$, take
$x^+=\xhat$. Otherwise choose a nearest point $f\in F$ and put
$t=2q/\tau_F$. Along the coordinate segment $z+s w_0$, $0\leq s\leq t$,
integration gives
\[
 G(m(z+t w_0))=f+t v_0+e,
 \qquad \norm{e}\leq q+\frac{t\tau_F}{4}
 =\frac{3}{4}t\tau_F<t\tau_F.
\]
The relative ball with center $v_0$ and radius $\tau_F$ is contained
in $F$. Cone homogeneity and $F+F=F$ therefore imply the displayed
endpoint belongs to $F$. The input path has length at most $2t$, proving
\eqref{eq:relative}. By continuity, $q\to0$ as $\xhat\to\xb$;
nesting neighborhoods ensures all these paths stay in the chart.
\end{proof}

\begin{theorem}[MSCQ from amenability of the reference face]
\label{thm:mscq}
Let $C$ be a proper nice cone, let $G$ be $C^1$ with $G(\xb)=0$, and
suppose frozen minimal-face CRSC holds at $\xb$. Assume the reference face
$F$ is amenable in $C$, that is, for some $a_F>0$,
\[
 \dset{y}{F}\leq a_F\dset{y}{C}\qquad(y\in\spann F).
\]
Then the system has MSCQ at $\xb$. This is a property of $F$
relative to the original cone $C$, not merely amenability of $F$
viewed as an independent cone.
For $F\ne\{0\}$, let $a_F$ satisfy \eqref{eq:amenable}, let $b$ be as
in Lemma~\ref{lem:normal-correction}, and let $L_G$ be a local
Lipschitz constant of $G$ on the neighborhoods used for both corrections.
Then one valid local upper bound is
\begin{equation}
 \boxed{\kappa=b+c\,a_F(1+L_Gb),\qquad c=4/\tau_F.}
 \label{eq:kappa}
\end{equation}
If $F=\{0\}$ and $b>0$, $\kappa=b$ is valid. If $F=\{0\}$ and
$b=0$, all sufficiently near inputs are feasible, and any $\kappa>0$
is valid.
\end{theorem}
\begin{proof}
Set $\rho=\dset{G(x)}{C}$. Lemma~\ref{lem:normal-correction} gives
$\xhat\in M$ with $\norm{x-\xhat}\leq b\rho$.
If $F=\{0\}$, then $\spann F=\{0\}$, so $G(\xhat)=0$ and
$\xhat$ is feasible. This proves the zero-face assertions.

Suppose $F\ne\{0\}$. Since $G(\xhat)\in\spann F$, amenability
and local Lipschitz continuity give
\begin{align}
 \dset{G(\xhat)}{F}
 &\leq a_F\dset{G(\xhat)}{C}\notag\\
 &\leq a_F\bigl(\rho+L_G\norm{x-\xhat}\bigr)
 \leq a_F(1+L_Gb)\rho.
 \label{eq:amenable-chain}
\end{align}
Apply Lemma~\ref{lem:relative} to this same corrected input. The resulting
$x^+$ belongs to $\Omega$, and
\[
 \dset{x}{\Omega}\leq\norm{x-x^+}
 \leq\norm{x-\xhat}+\norm{\xhat-x^+}
 \leq\bigl[b+c\,a_F(1+L_Gb)\bigr]\rho.
\]
Here the neighborhood is common to all inputs under consideration.
Fix first a product chart with the derivative bounds used in both lemmas,
then a smaller chart whose closure lies in it. The latter has positive
distance from the boundary of the larger coordinate domain. The normal
correction $(u,z)\mapsto(0,z)$ stays inside the product chart. Finally,
shrink one input neighborhood so that every corrected point is in the
smaller chart and every relative-correction coordinate segment has length
less than that fixed boundary margin. This is possible because
$\xhat\to\xb$ and $\dset{G(\xhat)}{F}\to0$. A convex input ball
containing all these points supplies the same Lipschitz constant $L_G$.
Only finitely many neighborhood restrictions are used, so the estimate
holds on one open neighborhood of $\xb$.
\end{proof}

\begin{corollary}[Amenable cones]
\label{cor:amenable}
For every proper amenable cone, frozen minimal-face CRSC at an apex
implies MSCQ for arbitrary $C^1$ constraint data.
\end{corollary}
\begin{proof}
Amenability implies niceness and supplies \eqref{eq:amenable} for the
reference face, so Theorem~\ref{thm:mscq} applies.
\end{proof}

\begin{remark}[Scope of the constants]
The constants in \eqref{eq:b}, \eqref{eq:tau}, and \eqref{eq:kappa}
are sufficient local upper bounds, not optimal subregularity moduli.
The proof establishes the existence of a sufficiently small radius;
it does not compute a universal radius from these constants alone.
For $F=C$ the rank-zero normal branch applies with $b=0$ and $a_F=1$,
leaving only the usual interior-direction correction.
The theorem requires amenability only of $F$ in $C$, not of all faces of $C$.
For reference, the degenerate cases are as follows:
\begin{center}\small
\begin{tabular}{@{}ll@{}}
\toprule
Case & Correction and bound\\
\midrule
$r=0$ & $\phi\equiv0$ locally; identity chart and $b=0$.\\
$F=\{0\}$ & Normal correction is already feasible; no $d_0$ or $\tau_F$.\\
$F=C$ & $H=\{0\}$, $b=0$, $a_F=1$; only interior correction.\\
$E=\{0\}$ reduction & Locally all feasible; any $\kappa>0$.\\
\bottomrule
\end{tabular}
\end{center}
In every case a zero face residual requires no second movement.
\end{remark}

\subsection{A universal characterization within nice cones}

\begin{proposition}[The isometric face-inclusion test]
\label{prop:inclusion}
Let $C$ be a closed convex cone in a finite-dimensional Euclidean space,
let $J$ be a closed face, and let $I_J:L\hookrightarrow E$ be inclusion,
where $L=\spann J$ has the induced Euclidean norm. Then
$I_J^{-1}(C)=J$, and this system has MSCQ at zero if and only if $J$ is
amenable in $C$. If $C$ is nice, the inclusion satisfies full facial
CRCQ at zero. Under the same niceness assumption, A1 and A2 hold globally
for this fixed apex representation.
\end{proposition}
\begin{proof}
First $C\cap\spann J=J$. Indeed, write $y\in C\cap\spann J$ as
$y=j_1-j_2$, with $j_1,j_2\in J$. Since $y+j_2=j_1\in J$ and $J$
is a face of a cone, $y\in J$. Inclusion is isometric, so its distance
to feasibility is exactly $\dset{y}{J}$ in the ambient metric.
Amenability therefore implies local MSCQ. Conversely, if some
$\varepsilon,\kappa>0$ give the MSCQ inequality for $y\in L$ with
$\norm{y}<\varepsilon$, apply it to $\lambda y$, with
$\lambda\norm{y}<\varepsilon$. Both cone distances scale by $\lambda$;
cancelling it proves the face estimate for every nonzero $y\in L$.
The estimate at zero is immediate.

Suppose now that $C$ is nice. Since $I_J^*=P_L$,
\begin{equation}
 I_J^*C^*=P_LC^*=(C^*+L^\perp)\cap L
 \label{eq:inclusion-image}
\end{equation}
is closed by niceness. For every face $Q$ of $C$, the image
$DI_J(y)^*[Q^\perp]=P_L[Q^\perp]$ is independent of $y$.
Thus all face-rank tests hold on one common neighborhood, indeed on
all of $L$, and reference closedness proves full facial CRCQ.
Finally $\range DI_J\cap C=J$ and $P_{J^\perp}DI_J=0$ globally.
The adjoint cone image and minimal linearized face are consequently
constant throughout this fixed representation.
\end{proof}

\begin{theorem}[Universal CRSC-to-MSCQ property]
\label{thm:universal}
Fix a proper nice cone $C$. The following are equivalent:
\begin{enumerate}[label=(\roman*)]
 \item $C$ is amenable.
 \item For every finite-dimensional Euclidean input space $X$, every open
 set $V\subset X$, every $\xb\in V$, and every $C^1$ map $G:V\to E$
 with $G(\xb)=0$, frozen minimal-face CRSC at $\xb$ implies the existence
 of $\kappa,\varepsilon>0$ such that $B(\xb,\varepsilon)\subset V$ and
 \[
 \dset{x}{G^{-1}(C)}\leq\kappa\dset{G(x)}{C}
 \qquad(x\in B(\xb,\varepsilon)).
 \]
 \item The statement in (ii) holds with full facial CRCQ in place of
 frozen minimal-face CRSC.
 \item For each face $J$ of $C$, the inclusion map
 $I_J:\spann J\hookrightarrow E$ has MSCQ at zero.
\end{enumerate}
The constants and radii may depend on the map and reference point, or on
the face in (iv). No common modulus or common radius across these
systems is asserted.
\end{theorem}
\begin{proof}
Corollary~\ref{cor:amenable} proves (i)$\Rightarrow$(ii).
Full facial CRCQ implies minimal-face CRSC, giving
(ii)$\Rightarrow$(iii). By Proposition~\ref{prop:inclusion}, niceness
makes every face inclusion satisfy full facial CRCQ, so
(iii)$\Rightarrow$(iv). The same proposition turns (iv) into amenability
of every face, proving (iv)$\Rightarrow$(i).
\end{proof}

The quantifiers are essential: this theorem characterizes a cone through
the behavior of \emph{all} admissible maps. It does not assert that
CRSC and MSCQ are equivalent for an individual map, nor that an
individual MSCQ system requires its entire output cone to be amenable.
The zero face has a zero-dimensional, locally feasible inclusion; the
full face has the identity inclusion and bound one. These cases require
no uniform convention across different input spaces. The short reverse
direction is consistent with the known facewise metric characterizations
in \citet[Proposition~12]{lourencoAmenable} and
\citet[Proposition~3.2]{lrs2022}. Niceness ensures the testing maps belong
to the quantified CQ class; the theorem is not asserted for arbitrary
closed cones without that standing assumption.

\section{The boundary: a nice-cone counterexample with constant derivative}
\label{sec:boundary}

For every nonamenable face of a nice cone,
Proposition~\ref{prop:inclusion} gives a linear full-facial-CRCQ system
that fails MSCQ. This does not mean every representation over a
nonamenable cone fails the error bound. It shows that even all frozen
facial rank tests cannot supply a missing face-to-cone residual estimate.

We next give an explicit witness using the cone of
\citet[Section~5]{lrs2022}. The cone, its niceness, and its original
tangency geometry are imported from that work. The presentation below
adapts them to an isometric constraint representation and a sequence
converging to the apex.

For $0\leq t\leq2\pi$, set
\[
 \alpha(t)=(\cos t,\sin t,1),\qquad
 \beta(t)=(\cos t,\sin t,-1).
\]
For $0\leq t\leq\pi$, set
\[
 \gamma(t)=\left(2\cos(2t)-1,\;2\sin(2t),\;
 \frac98\cos t-\frac18\cos(3t)\right).
\]
Here $\alpha,\beta,\gamma$ also denote the images of the three curves.
Define
\begin{equation}
 B=\conv(\alpha\cup\beta\cup\gamma),\qquad
 K=\cone(B\times\{1\})\subset\R^4.
 \label{eq:lrs-cone}
\end{equation}
The cone's niceness is \citet[Proposition~5.1]{lrs2022}; its
nonamenability and the existence statement appear in their
Proposition~5.2 and Theorem~5.3. Its upper disk generates the face
\begin{equation}
 J=\{(a,b,s,s):s\geq\sqrt{a^2+b^2}\}.
 \label{eq:lrs-face}
\end{equation}
On the Euclidean input space $\R^3$, take the linear isometry
\begin{equation}
 G(a,b,c)=(a,b,c/\sqrt2,c/\sqrt2).
 \label{eq:lrs-map}
\end{equation}
\begin{proposition}[Full facial CRCQ without MSCQ]
\label{prop:counterexample}
The cone $K$ in \eqref{eq:lrs-cone} is proper and nice, and $J$ in
\eqref{eq:lrs-face} is an exposed face. The system
$G(x)\in K$ with \eqref{eq:lrs-map} satisfies full facial CRCQ at zero,
but fails MSCQ there. Along an explicit path its distance to the whole
feasible set is $\Theta(t^3)$, while its actual positive cone residual
is $O(t^5)$.
\end{proposition}
\begin{proof}
The set $B$ is compact and contains the solid cylinder
$\conv(\alpha\cup\beta)$, so it has nonempty interior in $\R^3$.
Its height-one conic lift $K$ is full-dimensional and pointed. To see
closedness, if $\lambda_j(b_j,1)$ converges, its fourth coordinate
makes $\lambda_j$ converge. Compactness of $B$ treats a positive limit;
boundedness of $B$ forces the whole vector to zero if $\lambda_j\to0$.
Niceness is imported from the cited proposition.

Writing $c=\cos t$ gives
\[
 \gamma_3(t)=\frac32c-\frac12c^3,
 \qquad 1-\gamma_3(t)=\frac12(1-c)^2(2+c).
\]
On $[-1,1]$ the height is at most one, with equality only at $c=1$.
Hence $B\cap\{z=1\}=\conv\alpha$, and the nonnegative functional
$(a,b,z,s)\mapsto s-z$ exposes precisely $J$.
In particular $\range G=\spann J$, $K\cap\range G=J$, and the full
feasible set is
\begin{equation}
 \Omega=\{(a,b,c):c\geq\sqrt2\sqrt{a^2+b^2}\}.
 \label{eq:lrs-feasible}
\end{equation}
The isometry $G$ identifies $\R^3$ with $\spann J$.
Proposition~\ref{prop:inclusion} proves full facial CRCQ, including
reference adjoint-image closedness and one common neighborhood for all
face-rank tests.

For $0<t<\pi/2$, write
\[
 a(t)=2\cos(2t)-1,\quad b(t)=2\sin(2t),\quad
 r(t)=\sqrt{a(t)^2+b(t)^2}=\sqrt{5-4\cos(2t)},
\]
and choose inputs
$u(t)=t(a(t),b(t),\sqrt2)\to0$. Then
$G(u(t))=t(a(t),b(t),1,1)$.
Since $r(t)>1$, the output is outside $J$ but inside $\spann J$.
It is therefore outside the closed cone $K$, and the actual residual
$\delta_t=\dset{G(u(t))}{K}$ is positive.
Since $(\gamma(t),1)\in K$, homogeneity gives
\begin{align}
 0<\delta_t
 &\leq q_t:=t\left(1-\frac98\cos t+\frac18\cos(3t)\right)\notag\\
 &=\frac{t}{2}(1-\cos t)^2(2+\cos t)
 =\frac38t^5+O(t^7).
 \label{eq:small-residual}
\end{align}
This is an upper bound on the cone residual, not an asserted exact
distance to $K$.

Because $G$ is isometric and $K\cap\spann J=J$,
\[
 \dset{u(t)}{G^{-1}(K)}=\dset{G(u(t))}{J}.
\]
For the unscaled output, fixing $s\geq0$ and minimizing the first two
coordinates over the disk of radius $s$ gives the squared distance
\[
 \bigl(\max\{r(t)-s,0\}\bigr)^2+2(1-s)^2.
\]
On $0\leq s\leq r(t)$ this function is strictly convex with minimizer
$s_t=(r(t)+2)/3\in(1,r(t))$. On $s\geq r(t)$ it is minimized at
$r(t)$, which cannot improve that value. The global minimum is thus
$\frac23(r(t)-1)^2$. Isometry and cone homogeneity give
\begin{align}
 \dset{u(t)}{G^{-1}(K)}
 &=t\sqrt{\frac23}\,(r(t)-1)\notag\\
 &=4\sqrt{\frac23}\,t^3+O(t^5).
 \label{eq:large-distance}
\end{align}
Here $r(t)-1=8\sin^2t/(\sqrt{1+8\sin^2t}+1)$.
Consequently $\delta_t=O(t^5)$ and
\[
 \frac{\dset{u(t)}{\Omega}}{\delta_t}
 \geq\frac{\dset{u(t)}{\Omega}}{q_t},\qquad
 \lim_{t\downarrow0}t^2\frac{\dset{u(t)}{\Omega}}{q_t}
 =\frac{32}{3}\sqrt{\frac23}>0.
\]
For an explicit neighborhood contradiction, $0<t<\pi/2$ gives
$\sin t\geq2t/\pi$, $r(t)+1\leq4$, and
$1-\cos t\leq t^2/2$. Thus
\[
 \dset{u(t)}{\Omega}\geq\frac8{\pi^2}\sqrt{\frac23}\,t^3,
 \quad q_t\leq\frac38t^5,
 \quad \norm{u(t)}\leq\sqrt{11}\,t.
\]
Given any proposed $\varepsilon,\kappa>0$, set
$c_0=64\sqrt{2/3}/(3\pi^2)$ and choose
\[
 0<t<\min\left\{\frac\pi2,\frac{\varepsilon}{\sqrt{11}},
                  \sqrt{\frac{c_0}{\kappa}}\right\}.
\]
Then $\norm{u(t)}<\varepsilon$ while its actual distance-to-residual
ratio exceeds $\kappa$, contradicting MSCQ.
\end{proof}

The nearby face and closed-image properties also hold globally for this
linear map. Hence the obstruction is not instability of the rank or
minimal face. It is the failure of \eqref{eq:amenable} for $J$ inside $K$.
No new nice nonamenable cone is claimed. Only $q_t$, the constructed
upper bound, is shown here to be $\Theta(t^5)$; no matching lower bound
for $\delta_t$ or optimal H\"older exponent is asserted. Numerical
lower-bound-ratio checks for $\dset{u(t)}{\Omega}/q_t$ are merely
sanity checks of this analytic expression, not evaluations of the true
cone-distance ratio or certified interval computations.

\section{Cone classes, coupled examples, and fixed reductions}
\label{sec:examples}

\subsection{Products and strictly convex factors}

Let $C=\prod_{i=1}^m C_i$ in an orthogonal product space. Faces of a
product have the form $F=\prod_i F_i$. If $F_i$ is amenable in $C_i$
with constant $a_i$, then for $y\in\spann F$,
\[
 \dset{y}{F}^2=\sum_i\dset{y_i}{F_i}^2
 \leq(\max_i a_i)^2\sum_i\dset{y_i}{C_i}^2.
\]
Thus $a_F=\max_i a_i$ is valid. This is the established product
amenability calculus \citep{lourencoAmenable}, displayed here to fix
the constant and the output metric.

For a ray face $F_i=\R_+v_i$ with $\norm{v_i}=1$, the bound
\[
 a_i=\frac{1}{\dset{-v_i}{C_i}}
\]
is immediate by considering positive and negative multiples of $v_i$.
The denominator is positive by pointedness and closedness. Zero and full
faces have the trivial comparison. This proves the face estimates for
proper strictly convex cones, whose nonzero proper faces are rays.
In particular, finite products of the standard $p$-cones
\[
 C_p=\{(s,z):s\geq\norm{z}_p\},\qquad 1<p<\infty,
\]
and nonnegative scalar blocks are covered by
Corollary~\ref{cor:amenable}. The example classes and their amenability
are known; the conclusion used here is the joint nonlinear CRSC-to-MSCQ
implication.

This product calculation makes no assumption about blockwise MSCQ of
the input maps. A single $G$ may couple every variable across the blocks.
The proof first uses the one normal map $P_HG$, one rank, one angle,
and one common manifold, and only afterwards the amenability of the
whole output face.

\subsection{Positive-semidefinite faces}

Equip $\Ssym^m$ with the trace inner product. Every face of $\Ssym_+^m$
has the form
\[
 F=\{URU^{\mathsf T}:R\in\Ssym_+^k\},
 \qquad U^{\mathsf T}U=I_k.
\]
A matrix in $\spann F$ is $Y=UBU^{\mathsf T}$ with $B\in\Ssym^k$.
Its orthogonal projection onto the PSD cone is
$UB_+U^{\mathsf T}$, which already lies in $F$. Therefore
\begin{equation}
 \dset{Y}{F}=\dset{Y}{\Ssym_+^m}\qquad(Y\in\spann F),
 \label{eq:psd-face}
\end{equation}
and one may take $a_F=1$. The PSD specialization consequently includes
proper faces of arbitrary rank and dimension, not only ray faces.

\subsection{Two SOC blocks with one curved common equation}

The two examples below show joint structural coverage and the sources
of the proof constants, which we compare with direct estimates on a
specified ball.

Let $Q_3=\{(z_0,z_1,z_2):z_0\geq\sqrt{z_1^2+z_2^2}\}$, and put
\[
 v=(1,1,0),\qquad n=(1,-1,0),\qquad e=(0,0,1),
 \qquad h(s,t,u)=u-st.
\]
Define $G:\R^3\to\R^3\times\R^3$ by
\begin{equation}
 G_1=sv+hn+h^2e,\qquad
 G_2=tv-hn+2h^2e,
 \label{eq:soc-example}
\end{equation}
with feasibility requirement $G(s,t,u)\in Q_3\times Q_3$.
At the origin the derivative image is
\[
 (s,t,u)\longmapsto(sv+un,\;tv-un).
\]
The two pairings with $n$ are $2u$ and $-2u$, so membership in both cones
forces $u=0$, $s\geq0$, and $t\geq0$. The minimal face is thus
\[
 F=\R_+v\times\R_+v,
\]
a two-dimensional proper face.

The joint normal map is
\[
 P_HG=(hn+h^2e,\;-hn+2h^2e).
\]
Its derivative factors through $Dh=(-t,-s,1)$, which never vanishes.
The derivative with respect to $h$ is $(n+2he,-n+4he)$ with squared
norm $4+20h^2>0$, so the joint normal rank is one everywhere.
Reference closedness can be checked without another appeal to the
closed-image theorem:
\begin{equation}
 A^*C^*=\R_+^2\times\R.
 \label{eq:soc-adjoint}
\end{equation}
Indeed its first two coordinates are the nonnegative pairings with $v$.
Conversely, for $a,b\geq0$ and $z\in\R$, write
$z_+=\max\{z,0\}$, $z_-=\max\{-z,0\}$ and choose
\[
 y_1=(a/2)v+(z_+/2)n,\qquad y_2=(b/2)v+(z_-/2)n.
\]
Both belong to $Q_3$, and $A^*(y_1,y_2)=(a,b,z)$.
The conjugate face is $\R_+n\times\R_+n$, whose adjoint image, like
$A^*H$, is the full $u$-axis. Thus CRSC holds.

For the nonlinear map, the same pairings give $2h$ and $-2h$. Feasibility
forces $h=0$, and conversely this removes all normal terms. Hence
\begin{equation}
 \Omega=\{(s,t,u):s\geq0,\ t\geq0,\ u=st\}.
 \label{eq:soc-feasible}
\end{equation}
The common face manifold is the curved surface $h=0$.
Each block separately has an interior derivative direction: for example,
$s=u>0$ for the first block gives a positive multiple of $(1,0,0)$,
and $t=-u>0$ does the same for the second. The joint derivative has no
interior direction because its two normal pairings have opposite signs.
Theorem~\ref{thm:mscq} applies jointly to the nonisolated feasible set
\eqref{eq:soc-feasible}.

\paragraph{Geometric constants and a general local bound.}
Use the orthonormal vectors $\widehat v=v/\sqrt2$,
$\widehat n=n/\sqrt2$, and $e$. The normal projection closure is
\begin{equation}
 \cl(P_{v^\perp}Q_3)
 =\{a\widehat n+ze:a\geq0,\ z\in\R\}.
 \label{eq:soc-closure}
\end{equation}
To see why the closure is needed, write a full vector as
$\lambda\widehat v+a\widehat n+ze$. Its SOC conditions are
$\lambda,a\geq0$ and $2\lambda a\geq z^2$. When $a>0$, any $z$ is
possible by increasing $\lambda$; at $a=0$ the original projection
requires $z=0$, whereas the closure permits every $z$.
The two unit directions of $U=\spann\{(n,-n)\}$ are
$\pm(n,-n)/2$, each at distance $1/\sqrt2$ from the product of the
closed half-spaces in \eqref{eq:soc-closure}. Since the sole positive
singular value of $P_HA$ is $2$, the constants are
\[
 \sigma=2,\qquad\eta=1/\sqrt2,\qquad b=2\sqrt2.
\]
The unit tangent direction $d_0=(1,1,0)/\sqrt2$ has image
$(\widehat v,\widehat v)$, so $\tau_F=1$ and $c=4$.
On the line spanned by $\widehat v$, the SOC projection of
$-\widehat v$ is zero. The face comparison is consequently exact there;
orthogonal product distances give $a_F=1$.

For $\norm{x}\leq R$, set $H_R=R+R^2/2$. Then $|h|\leq H_R$ and
$\norm{Dh}^2\leq1+R^2$. The tangential output increment
$(\Delta s\,v,\Delta t\,v)$ is orthogonal to the normal increment, hence
\[
 \norm{DG(x)}^2\leq2+(4+20H_R^2)(1+R^2).
\]
At $R=1/4$ the right side is $32485/4096<9$, so $L_G=3$ is valid
on that ball. Substitution in \eqref{eq:kappa} yields
\begin{equation}
 \kappa_{\rm th}=4+26\sqrt2
 \label{eq:soc-kappa-th}
\end{equation}
on some sufficiently small neighborhood. The calculation of $L_G$ on
$B_{1/4}$ does not, by itself, identify the radius required by the
general constant-rank coordinate proof.

\paragraph{A direct bound on a specified ball.}
Write $s_- =\max\{-s,0\}$, $t_- =\max\{-t,0\}$,
$a^2=s_-^2+t_-^2$, and $\rho=\dset{G(x)}{C}$. Each SOC lies in the
intersection of the two orthogonal half-spaces defined by nonnegative
pairings with $\widehat v$ and $\widehat n$. The two block residuals
therefore give
\[
 \rho^2\geq2(s_-^2+t_-^2+h^2)=2(a^2+h^2).
\]
Set $x^+=(s_+,t_+,s_+t_+)\in\Omega$.
For $\norm{x}\leq R$, Cauchy--Schwarz gives
\[
 |st-s_+t_+|\leq |t|s_-+s_+t_-
 \leq\sqrt{t^2+s_+^2}\,a\leq Ra.
\]
Thus $\norm{x-x^+}\leq\sqrt{a^2+h^2}+Ra\leq(1+R)\rho/\sqrt2$,
which proves the independently quantified estimate
\begin{equation}
 \dset{x}{\Omega}\leq\frac5{4\sqrt2}\dset{G(x)}{C}
 \qquad(\norm{x}\leq1/4).
 \label{eq:soc-direct}
\end{equation}
The constructed point need not be a nearest feasible point, and neither
this constant nor \eqref{eq:soc-kappa-th} is claimed to be optimal.

\subsection{A PSD model with a high-dimensional nonpolyhedral face}

Let $k\geq3$. Use inputs $(B,u)\in\Ssym^k\times\R$ with norm
$\sqrt{\norm{B}_F^2+u^2}$, and the Frobenius norm on the output.
Set $h=u-\tr(B^2)=u-\norm{B}_F^2$, and fix
\[
 C_0=[e_1,e_2],\qquad D_0=[e_2,e_3],\qquad
 J_0=\diag(1,-1),
\]
where $e_i$ are the standard vectors of $\R^k$.
Thus $\norm{C_0}_F^2=\norm{D_0}_F^2=2$ and
$\ip{C_0}{D_0}=0$. Define
\begin{equation}
 G(B,u)=
 \begin{pmatrix}
 B & hC_0+h^2D_0\\
 hC_0^{\mathsf T}+h^2D_0^{\mathsf T} & hJ_0
 \end{pmatrix}\in\Ssym^{k+2}.
 \label{eq:psd-example}
\end{equation}
At zero, a PSD derivative value must have $u=0$ by its bottom principal
block. Its off-diagonal block then vanishes. Thus the minimal face is
\begin{equation}
 F=\left\{\begin{pmatrix}B&0\\0&0\end{pmatrix}:B\succeq0\right\},
 \qquad \dim F=\frac{k(k+1)}{2}.
 \label{eq:psd-minimal}
\end{equation}
This face is exposed by $\diag(0_k,I_2)$ and is isometric to
$\Ssym_+^k$. It is nonpolyhedral, since rank-one PSD matrices along
infinitely many directions give distinct extreme rays. The smallest
case here, $k=3$, has a six-dimensional proper face of $\Ssym_+^5$.

For a symmetric block matrix $Y=\left(\begin{smallmatrix}P&Z\\
Z^{\mathsf T}&W\end{smallmatrix}\right)$,
\[
 A^*Y=(P,\,2\ip{C_0}{Z}+\ip{J_0}{W}).
\]
If $Y\succeq0$ then $P\succeq0$. Conversely any
$(P,z)\in\Ssym_+^k\times\R$ is obtained with $Z=0$ and
$W=\diag(z_+,z_-)$. Hence
\begin{equation}
 A^*C^*=\Ssym_+^k\times\R,
 \qquad A^*F^\triangle=A^*H=\{0\}\times\R.
 \label{eq:psd-adjoint}
\end{equation}
In particular the adjoint cone image is closed.
Write $q(h)=P_HG$. Its derivative and squared Frobenius norm are
\[
 q'(h)=\begin{pmatrix}0&C_0+2hD_0\\
 C_0^{\mathsf T}+2hD_0^{\mathsf T}&J_0\end{pmatrix},
 \qquad\norm{q'(h)}_F^2=6+16h^2>0.
\]
Since $Dh(E,t)=t-2\ip{B}{E}$ always has $u$-component one,
$D(P_HG)=q'(h)Dh$ has rank one everywhere. This verifies the full
neighborhood condition for CRSC, using the joint normal map.

For the nonlinear system, the bottom principal block again forces $h=0$.
All cross terms then vanish, giving the exact equivalence
\begin{equation}
 G(B,u)\succeq0
 \quad\Longleftrightarrow\quad
 B\succeq0,\qquad u=\tr(B^2).
 \label{eq:psd-feasible}
\end{equation}
The theorem supplies MSCQ for this representation with $a_F=1$.
The example demonstrates coupled nonlinear normal components and
arbitrarily large nonpolyhedral faces; it does not claim that elementary
direct analysis could not establish an error bound for this example.

\paragraph{Projection closure and general local constants.}
In this example
\begin{equation}
 D=\cl(P_H\Ssym_+^{k+2})
 =\left\{\begin{pmatrix}0&Z\\Z^{\mathsf T}&W\end{pmatrix}:
             W\succeq0,\ Z\in\R^{k\times2}\right\}.
 \label{eq:psd-closure}
\end{equation}
The lower block of any PSD matrix proves one inclusion. For the other,
replace $W$ by $W+\varepsilon I_2$ and choose the upper block
$Z(W+\varepsilon I_2)^{-1}Z^{\mathsf T}$. The resulting full matrix
is PSD, and its normal projection converges to the specified point as
$\varepsilon\downarrow0$. Singular lower blocks would impose extra
cross-block restrictions before taking the closure.

Let $N_0=q'(0)$, so $\norm{N_0}_F^2=6$ and $P_HA(E,t)=tN_0$.
The cross block can be retained exactly in a projection onto $D$;
the distance comes only from projecting the lower block onto
$\Ssym_+^2$. Each of $J_0$ and $-J_0$ has one negative eigenvalue
of magnitude one. Consequently
\[
 \sigma=\sqrt6,\qquad
 \eta=\dset{\pm N_0/\sqrt6}{D}=1/\sqrt6,\qquad b=4.
\]
The normal cross terms affect both $\sigma$ and the normalized angle;
their effects cancel in $b$. The unit tangent direction
$d_0=(I_k/\sqrt k,0)$ gives $\tau_F=1/\sqrt k$, since the Frobenius
distance to the relative PSD boundary is the smallest eigenvalue.
Thus $c=4\sqrt k$ and, by \eqref{eq:psd-face}, $a_F=1$.

On $\norm{(B,u)}\leq R$, set $H_R=R+R^2$.
The bounds $|h|\leq H_R$, $\norm{Dh}^2\leq1+4R^2$, and orthogonality
of the upper-block and normal increments yield
\[
 \norm{DG(B,u)}^2\leq1+(6+16H_R^2)(1+4R^2).
\]
At $R=1/4$ the right side is $669/64<(13/4)^2$, so $L_G=13/4$ is
valid on this specified ball. The general theorem then gives
\begin{equation}
 \kappa_{\rm th}=4+56\sqrt k
 \label{eq:psd-kappa-th}
\end{equation}
on some sufficiently small neighborhood, not an automatically
computed radius of $1/4$.

\paragraph{Direct spectral repair on the specified ball.}
Decompose $B=B_+-B_-$ into orthogonal positive and negative spectral
parts. Put $a=\norm{B_-}_F$ and
$\rho=\dset{G(B,u)}{\Ssym_+^{k+2}}$.
The two principal blocks of any PSD comparison matrix are PSD; discarding
the nonnegative cross-block distance term gives
\[
 \rho^2\geq\dset{B}{\Ssym_+^k}^2
              +\dset{hJ_0}{\Ssym_+^2}^2=a^2+h^2.
\]
The point $(B^+,u^+)=(B_+,\norm{B_+}_F^2)$ is feasible, and
$u-u^+=h+a^2$. If $\norm{(B,u)}\leq R$, then $a\leq R$ and
\[
 \norm{(B,u)-(B^+,u^+)}
 =\sqrt{a^2+(h+a^2)^2}
 \leq\sqrt{a^2+h^2}+a^2\leq(1+R)\rho.
\]
Therefore the independent, explicit-radius estimate is
\begin{equation}
 \dset{(B,u)}{\Omega}\leq\frac54\dset{G(B,u)}{\Ssym_+^{k+2}}
 \qquad(\norm{(B,u)}\leq1/4).
 \label{eq:psd-direct}
\end{equation}
This dimension-independent direct bound shows that the $\sqrt k$
growth in \eqref{eq:psd-kappa-th} belongs to the selected general
direction and proof constants, not to a demonstrated necessary
dimension dependence of this example. No optimal modulus is claimed.

\begin{table}[t]
\centering\small
\caption{Constants for the two joint examples. The Lipschitz bounds are
on $B_{1/4}$; only the final error-bound row is proved directly on that whole ball.}
\label{tab:constants}
\begin{tabular}{@{}lll@{}}
\toprule
Quantity & Coupled SOC & PSD, $k\geq3$\\
\midrule
$\sigma$ & $2$ & $\sqrt6$\\
$\eta$ & $1/\sqrt2$ & $1/\sqrt6$\\
$a_F$ & $1$ & $1$\\
$\tau_F$ & $1$ & $1/\sqrt k$\\
$L_G$ on $B_{1/4}$ & $3$ & $13/4$\\
$\kappa_{\rm th}$, unspecified local radius & $4+26\sqrt2$ & $4+56\sqrt k$\\
Direct $\kappa$ on $B_{1/4}$ & $5/(4\sqrt2)$ & $5/4$\\
\bottomrule
\end{tabular}
\end{table}

\subsection{Transfer through a fixed local reduction}
\label{subsec:reduction}

\begin{lemma}[Output-local residual transfer]
\label{lem:transfer}
Let $K$ be a closed convex cone, let $g$ be continuous near $\xb$, and
put $\bar y=g(\xb)\in K$, not necessarily an apex point.
Suppose there is an open output neighborhood $O$ of
$\bar y$ and a fixed $C^1$ map $\Xi:O\to E$ such that
\begin{equation}
 \Xi(\bar y)=0,\qquad
 y\in K\ \Longleftrightarrow\ \Xi(y)\in C
 \quad(y\in O).
 \label{eq:reduction}
\end{equation}
Assume the reduced data $G=\Xi\circ g$ meet the hypotheses of
Theorem~\ref{thm:mscq}. Choose a positive local Lipschitz bound
$L_\Xi$ for $\Xi$. Then
\begin{equation}
 \dset{G(x)}{C}\leq L_\Xi\dset{g(x)}{K}
 \label{eq:residual-transfer}
\end{equation}
holds near $\xb$, and the original system has MSCQ with constant
$\kappa L_\Xi$, where $\kappa$ is a bound supplied by
Theorem~\ref{thm:mscq}.
\end{lemma}
\begin{proof}
Projection onto the closed convex cone $K$ is continuous. Since
$g(x)\to\bar y=P_K\bar y$, shrink the input neighborhood until both
$g(x)$ and $P_Kg(x)$ lie in one output neighborhood on which $\Xi$ has
the chosen Lipschitz bound. Then \eqref{eq:reduction} yields
\[
 \dset{G(x)}{C}
 \leq\norm{\Xi(g(x))-\Xi(P_Kg(x))}
 \leq L_\Xi\dset{g(x)}{K}.
\]
The nearby corrected feasible points in Theorem~\ref{thm:mscq} are
also feasible for the original system by \eqref{eq:reduction}.
Its distance estimate therefore proves the claimed original bound.
Only this residual direction is needed; neither a reverse comparison
nor surjectivity of $D\Xi$ is used in this transfer step.
\end{proof}

For PSD constraints, a standard fixed Schur-complement reduction is
explicit. In an orthogonal basis at a reference PSD matrix of rank $q$,
write a nearby symmetric matrix as
\[
 Y=\begin{pmatrix}Y_{11}&Y_{12}\\Y_{12}^{\mathsf T}&Y_{22}\end{pmatrix},
 \qquad Y_{11}\succ0,
\]
and take
$\Xi(Y)=Y_{22}-Y_{12}^{\mathsf T}Y_{11}^{-1}Y_{12}$.
Then $Y\succeq0$ if and only if $\Xi(Y)\succeq0$, and the reduced
reference is zero. If $q$ is the full matrix size, the reduced output
space is zero-dimensional and the separate locally feasible branch
applies. We do not infer such a local reduction at every nonapex point
of every amenable cone from amenability alone.

\section{A regularity and residual-transfer interface}
\label{sec:residual-interface}

The natural relation certified by Theorem~\ref{thm:mscq} is
\[
 \mathcal M(x)=G(x)-C,\qquad
 \dset{0}{\mathcal M(x)}=\dset{G(x)}{C}.
\]
Its regularity gauge is $\psi_{\mathcal M}(t)=\kappa t$.
For a set-valued operator $\mathcal T$, write
$Z_{\mathcal T}=\{x:0\in\mathcal T(x)\}$. Assume independently that
$Z_{\mathcal T}$ agrees with $\Omega$ locally and that
\begin{equation}
 \dset{G(x)}{C}\leq
 \chi\bigl(\dset{0}{\mathcal T(x)}\bigr)
 \label{eq:operator-comparison}
\end{equation}
holds near $\xb$, with $\chi$ nondecreasing and vanishing at zero.
The same nearby feasible corrections then give the local bound
\begin{equation}
 \dset{x}{Z_{\mathcal T}}
 \leq\kappa\chi\bigl(\dset{0}{\mathcal T(x)}\bigr),
 \qquad \psi_{\mathcal T}(t)=\kappa\chi(t).
 \label{eq:operator-gauge}
\end{equation}
This only transfers a regularity certificate. The relation
$\mathcal M:X\rightrightarrows E$ need not be a square-space operator
for a proximal-point iteration; the operator, its geometric hypotheses,
and the zero-set/residual comparisons remain separate obligations.
If $\chi$ is linear, the resulting gauge remains linear. No nonlinear
exponent or proximal-point convergence rate is asserted.

\section{Limitations and conclusion}
\label{sec:conclusion}

The positive metric result is deliberately limited to a proper nice
cone with an amenable reference face. Nonpointed and non-full-dimensional
models would require explicit restriction or quotient constructions and
are outside the frozen scope of this manuscript. Nonapex applications
require a fixed valid local reduction as in
Section~\ref{subsec:reduction}; no automatic reducibility assertion is
made for arbitrary amenable cones. All rank statements use the reference
face on an open input neighborhood. The result does not replace that
condition by rank constancy only along the feasible set.

The argument is analytic and uses $C^1$ data. It is not a numerical
verification procedure for exact neighborhood rank constancy, and the
displayed constants do not determine an optimal modulus or an explicit
largest validity radius. The illustrative SOC and PSD systems demonstrate
the scope of a joint map, not algorithmic superiority. The nice
nonamenable cone and the foundational cone-geometric results retain their
original attribution.

The principal conclusion is the separation of geometric and metric
roles. Frozen minimal-face CRSC over a proper nice cone already controls
the missing nearby face geometry. A quantitative correction to the
common face manifold then allows the face's amenability to be used with
the original residual. Combining this sufficient direction with linear
face inclusions gives an exact universal boundary inside the nice class:
amenability is precisely the property under which every apex CRSC system
has MSCQ. The same boundary explains why full facial rank constancy
cannot restore a linear error bound in the nonamenable example.

\bibliographystyle{plainnat}
\bibliography{references}

\end{document}
